\documentclass[12pt,a4paper]{article}

\usepackage{amsmath,amssymb}
\usepackage{hyperref}
\usepackage{natbib}
\usepackage{geometry}
\usepackage{booktabs}
\usepackage{xcolor}
\usepackage{microtype}

\geometry{margin=1in}

\hypersetup{
  colorlinks=true,
  linkcolor=blue,
  citecolor=blue,
  urlcolor=blue
}


\title{\textbf{Entanglement, Decoherence, and the Thermodynamic Cost\\
of Classical Records}}
\author{Heath W.\ Mahaffey\\
\small Independent Researcher, Entiat, Washington 98822, USA\\
\small \href{mailto:hmahaffeyges@gmail.com}{hmahaffeyges@gmail.com}\\
\small \url{https://github.com/hmahaffeyges/IAM-Validation}}
\date{March 2026; revised October 2026}

\begin{document}
\maketitle

\begin{abstract}
Entangled systems show correlations that violate Bell inequalities
\citep{Bell1964,Aspect1982,Hensen2015}, ruling out local hidden-variable models.
Quantum mechanics predicts these correlations exactly. What it leaves open is the
quantum-to-classical transition: where and at what cost a superposition becomes a
definite classical record. Decoherence theory \citep{Zurek1981,JoosZeh1985,Zurek2003}
explains the suppression of interference; Landauer's principle \citep{Landauer1961}
sets the minimum thermodynamic cost of an irreversible record, $k_BT\ln2$ per bit.
The Informational Actualization Model (IAM) applies that cost at cosmological scale:
at the Gibbons--Hawking temperature $T_H=\hbar H/2\pi k_B$ the informational entropy
produced by structure formation enters horizon thermodynamics, giving
$\mu(a)<1$ for matter and $\Sigma(a)=1$ for light \citep{IAM_theory}. Here we set out
what this picture says about entanglement. It makes no change to quantum
mechanics and adds no hidden variables: Bell correlations are those of the shared
quantum state until an irreversible record is made. Its new content is
quantitative: a gravitational decoherence time
$\tau_{\rm IAM}=\hbar k_B^2T^2\ln2/E_G^3$ that scales as $T^2$ and, at fixed
density, as $m^{-5}$, against $m^{-5/3}$ and no temperature dependence for
Penrose--Di\'osi. For a 1~pg silica sphere at 10~mK, $\tau_{\rm IAM}=559$~s and
$\tau_{\rm PD}=7.8~\mu$s; the two cross at $m\approx2.3\times10^{-10}$~kg. We also
correct the exponent that links the cosmological activation function to structure
formation: matching $\exp(1-1/a)$ in matter domination requires $n=7/2$.
\end{abstract}

% ──────────────────────────────────────────────────────────────────
\section{Entanglement and the quantum-to-classical transition}
\label{sec:intro}

Bell's theorem \citep{Bell1964} shows that no local hidden-variable theory reproduces
the correlations of entangled states. Experiments confirm the violation, including a
loophole-free test with electron spins 1.3~km apart \citep{Hensen2015} and photon
pairs distributed over 1,200~km by satellite \citep{Yin2017}. Quantum mechanics
predicts these results without modification.

Two distinct questions are often merged. The first is why the correlations are
non-classical; the answer is that an entangled pair is described by one quantum
state, not by two systems with pre-existing local properties. The second is where
and how that state gives way to definite classical outcomes. Decoherence explains
the loss of interference through entanglement with an environment
\citep{Zurek1981,JoosZeh1985}, and quantum Darwinism explains how redundant records
make outcomes objective \citep{Zurek2003,Zurek2009}. Neither fixes the
thermodynamic cost of the record. Landauer's principle does: an irreversible
record costs at least $k_BT\ln2$ per bit, dissipated as heat \citep{Landauer1961}.

IAM takes this cost seriously at the largest scale. If structure formation produces
classical information irreversibly, the corresponding entropy belongs in the
horizon thermodynamics from which the Friedmann equations follow
\citep{Jacobson1995,CaiKim2005}. This paper states what that does and does not
imply for entanglement.

% ──────────────────────────────────────────────────────────────────
\section{Two sectors}
\label{sec:sectors}

Gravitational decoherence in IAM requires a timelike worldline, $d\tau>0$. Photons
travel null geodesics, $d\tau=0$, and undergo no gravitational decoherence in
transit; they are absorbed and measured in the ordinary way, but propagation adds
nothing to the informational entropy. Massive particles accumulate proper time and,
when bound into virialised structures, contribute to it. This is the origin of the
matter/photon sector split \citep{IAM_time}.

Local environmental decoherence (gas collisions, thermal photons, phonons) is the
dominant mechanism in laboratory systems and is unchanged by IAM. IAM adds a
gravitational channel and an accounting statement: every irreversible record made
on a timelike worldline carries its Landauer cost, and at cosmological scale that
cost appears in the horizon entropy budget.

% ──────────────────────────────────────────────────────────────────
\section{The derivation chain}
\label{sec:derivation}

\textbf{Step 1.} Jacobson derived the Einstein equations as an equation of state from
the entropy--area relation of local causal horizons \citep{Jacobson1995}.

\textbf{Step 2.} Cai and Kim obtained the Friedmann equations from the first law on
the apparent horizon of an FRW universe \citep{CaiKim2005}.

\textbf{Step 3.} Structure formation produces classical information irreversibly.
For $1/r$ potentials the virial theorem gives $2K=|V|$. IAM identifies the kinetic
half with the heat dissipated in the quantum-to-classical transition, priced at
$k_BT_H\ln2$ per bit \citep{Gibbons1977}.

\textbf{Step 4.} The accumulated informational contribution is described by the
activation function
\begin{equation}
\mathcal{E}(a)=\exp\!\left(1-\frac{1}{a}\right),
\label{eq:Ea}
\end{equation}
monotonically increasing, $\mathcal{E}(1)=1$, $\mathcal{E}(\infty)=e$, and
negligible at early times ($\mathcal{E}(z=10)=4.5\times10^{-5}$).

\textbf{Step 5.} Including this entropy in the first law adds Hubble friction to the
matter perturbation equation:
\begin{equation}
\mu(a)=\frac{H^2_{\Lambda\rm CDM}}{H^2_{\Lambda\rm CDM}+\beta_m\mathcal{E}(a)}<1,
\qquad \Sigma(a)=1,
\label{eq:sector}
\end{equation}
with $\beta_m=\Omega_m/2=0.15765$ from the virial theorem ($\Omega_m=0.3153$). The
Planck 2018 Level~2 posterior returns $\beta_m=0.1583\pm0.0033$, $0.2\sigma$ from
this value \citep{IAM_obs}.

% ──────────────────────────────────────────────────────────────────
\section{Entanglement under the two sectors}
\label{sec:entanglement}

Because IAM leaves quantum mechanics unchanged, Bell correlations of a pure shared
state are those of standard theory. The framework's statement is about how long an
entangled state of \emph{massive} systems can survive the gravitational channel.

For photon pairs no gravitational decoherence acts in transit, and the CHSH
parameter stays at $|S|=2\sqrt2$ for an ideal source, whatever the distance. For
massive superpositions IAM predicts gravitational decoherence with the time-dependent
profile of \citet{IAM_QM},
\begin{equation}
S(t)=2\sqrt2\,\bigl(1-D(t)\bigr),\qquad D(t)=\mathcal{E}(t/\tau_{\rm IAM})/e ,
\label{eq:bell}
\end{equation}
and timescale \citep{IAM_grav}
\begin{equation}
\tau_{\rm IAM}=\frac{\hbar k_B^2T^2\ln2}{E_G^3},\qquad E_G=\frac{Gm^2}{r},
\label{eq:tauIAM}
\end{equation}
where $T$ is the bath temperature. At fixed density, $E_G\propto m^{5/3}$, so
$\tau_{\rm IAM}\propto T^2m^{-5}$; the Penrose--Di\'osi time
$\tau_{\rm PD}=\hbar/E_G\propto m^{-5/3}$ has no temperature dependence.

\begin{table}[h]
\centering
\small
\caption{Gravitational decoherence times for silica spheres
($\rho=2000$~kg\,m$^{-3}$), computed from Eq.~(\ref{eq:tauIAM}) and
$\tau_{\rm PD}=\hbar/E_G$.}
\label{tab:decoherence}
\begin{tabular}{lccc}
\toprule
System & $\tau_{\rm IAM}$ & $\tau_{\rm PD}$ & Discriminator \\
\midrule
1 fg sphere ($10^{-15}$ kg), 10 mK & $5.6\times10^{17}$ s & 0.78 s & PD testable; IAM not \\
1 pg sphere ($10^{-12}$ kg), 10 mK & 559 s & $7.8~\mu$s & $T^2$ vs $T^0$; ramp vs exponential \\
1 pg sphere, 20 mK & 2,238 s & $7.8~\mu$s & doubling $T$ quadruples $\tau_{\rm IAM}$ \\
Crossover, 10 mK & \multicolumn{2}{c}{$\tau_{\rm IAM}=\tau_{\rm PD}$ at $m\approx2.3\times10^{-10}$ kg} & \\
\bottomrule
\end{tabular}
\end{table}

The discriminating experiment is therefore a mesoscopic superposition near
$10^{-12}$~kg: Penrose--Di\'osi predicts loss of coherence in microseconds, IAM
predicts none on that timescale, and a temperature scan separates the two. Distance
records do not discriminate: matter-qubit entanglement over 1.3~km
\citep{Hensen2015} and photon entanglement over 1,200~km \citep{Yin2017} are both
consistent with standard decoherence and photon-mediated distribution.

Entanglement of top-quark pairs at the LHC \citep{CMS2024} is consistent with IAM:
the top lifetime, $\sim5\times10^{-25}$~s, is far shorter than any gravitational
decoherence time. Bell violations from path indistinguishability without
conventional entanglement \citep{SciAdv2025} concern the structure of the state and
are outside the gravitational channel. Correlations between photons that never
coexisted \citep{Megidish2013} are properties of the quantum state across time; a
treatment within IAM is left open.

% ──────────────────────────────────────────────────────────────────
\section{The exponent linking structure formation to $\mathcal{E}(a)$}
\label{sec:exponent}

Write the information production rate during matter domination as
$\dot I\propto\rho_m D^n H$, with $\rho_m\propto a^{-3}$, $D\propto a$ and
$H\propto a^{-3/2}$. With the Gibbons--Hawking temperature and the horizon area,
the entropy increment per unit $a$ scales as $a^{\,n-11/2}$, so that
\begin{equation}
S_{\rm info}(a)\propto\int a^{\,n-11/2}\,da\propto a^{\,n-9/2}.
\end{equation}
Matching the $-1/a$ dependence of $\ln\mathcal{E}(a)$ requires
\begin{equation}
n-\frac{9}{2}=-1\quad\Longrightarrow\quad n=\frac{7}{2}.
\label{eq:n}
\end{equation}
(For $n<9/2$ the integral is dominated by its lower limit, so $S_{\rm info}$ is
defined up to a constant fixed by the normalisation $\mathcal{E}(1)=1$.) A
bottom-up value of $n$ from the halo mass function must reproduce $7/2$ for the
chain to close; the numerical halo-collapse calculation gives an effective $1/a$
coefficient near unity at $\sigma_*\approx1.5$ \citep{IAM_Zurek}. Deriving $n$
analytically from the mass function is an open problem.

% ──────────────────────────────────────────────────────────────────
\section{Observational status}
\label{sec:obs}

Seventeen MCMC chains (12 MGCAMB, 5 modified CAMB) were run with the Planck 2018
likelihood (TT, TE, EE + low-$\ell$ + lensing), some combined with BAO, Pantheon+ or
RSD data; final Gelman--Rubin $R-1$ is below 0.01 for 15 and 0.023--0.024 for the
other two \citep{IAM_obs}. Results:
\begin{itemize}
\item $\sigma_8=0.7998\pm0.0058$ (ΛCDM $0.809$);
\item $H_0^{(\rm photon)}=67.16\pm0.47$ and $H_0^{(\rm matter)}=72.26$ km\,s$^{-1}$\,Mpc$^{-1}$;
\item $\beta_m=0.1583\pm0.0033$, $0.2\sigma$ from $\Omega_m/2=0.15765$;
\item $\Delta\chi^2=+0.54$ at best fit relative to ΛCDM.
\end{itemize}
Planck is consistent with the IAM modification but does not by itself establish it.
The prediction $\mu_0-1=-0.136$ is tested by Euclid ($\sigma(\mu_0)\approx0.04$,
$3.4\sigma$) and Euclid with DESI ($5.4\sigma$); $\Sigma_0=0$ is tested by Euclid
weak lensing.

% ──────────────────────────────────────────────────────────────────
\section{IAM's law}
\label{sec:law}

\noindent\fbox{\parbox{0.93\textwidth}{
\textbf{IAM's law.} Every irreversible quantum-to-classical transition on a timelike
worldline produces a classical record at a thermodynamic cost of at least
$k_BT\ln2$ per bit, with $T$ the temperature of the system that absorbs it; at
cosmological scale, $T=T_H$ and the cost enters the horizon entropy.
}}

\medskip
Consequences: (1) quantum mechanics is unchanged, and Bell correlations are those of
the shared state; (2) gravitational decoherence of massive superpositions follows
Eq.~(\ref{eq:tauIAM}), with a $T^2$ signature; (3) $\mathcal{E}(a)$ increases
monotonically because each record is irreversible.

% ──────────────────────────────────────────────────────────────────
\section{Discussion}
\label{sec:discussion}

Decoherence has been known to be irreversible since the 1980s
\citep{Zurek1981,JoosZeh1985}; horizon entropy since Bekenstein and Hawking
\citep{Bekenstein1973,Hawking1975}; the thermodynamic origin of gravity since
Jacobson \citep{Jacobson1995}. IAM asks where the Landauer cost of cosmological
structure formation goes, and answers that it enters the horizon entropy budget,
producing $\mu<1$ and $\Sigma=1$. For entanglement the framework predicts nothing
beyond quantum mechanics for photons, and a specific, temperature-dependent
gravitational decoherence time for massive superpositions that mesoscopic
experiments can test.

\section*{Acknowledgments}
The thermodynamic framework of T.\ Jacobson \citep{Jacobson1995} and the quantum
Darwinism programme of W.\,H.\ Zurek \citep{Zurek2003,Zurek2009} are the foundations
of this work.

\section*{Data Availability}
Chains, scripts and figures: \url{https://github.com/hmahaffeyges/IAM-Validation};
Zenodo DOI \href{https://doi.org/10.5281/zenodo.18702042}{10.5281/zenodo.18702042}.

\begin{thebibliography}{99}
\bibitem[Aspect et al.(1982)]{Aspect1982} Aspect, A., Dalibard, J., \& Roger, G.\ 1982, Phys.\ Rev.\ Lett., 49, 1804
\bibitem[Bekenstein(1973)]{Bekenstein1973} Bekenstein, J.\ D.\ 1973, Phys.\ Rev.\ D, 7, 2333
\bibitem[Bell(1964)]{Bell1964} Bell, J.\ S.\ 1964, Physics, 1, 195
\bibitem[Cai \& Kim(2005)]{CaiKim2005} Cai, R.-G.\ \& Kim, S.\ P.\ 2005, JHEP, 02, 050
\bibitem[CMS Collaboration(2024)]{CMS2024} CMS Collaboration 2024, Observation of quantum entanglement in top quark pair production, Rep.\ Prog.\ Phys., 87, 117801
\bibitem[Gibbons \& Hawking(1977)]{Gibbons1977} Gibbons, G.\ W.\ \& Hawking, S.\ W.\ 1977, Phys.\ Rev.\ D, 15, 2738
\bibitem[Hawking(1975)]{Hawking1975} Hawking, S.\ W.\ 1975, Commun.\ Math.\ Phys., 43, 199
\bibitem[Hensen et al.(2015)]{Hensen2015} Hensen, B., et al.\ 2015, Nature, 526, 682
\bibitem[Jacobson(1995)]{Jacobson1995} Jacobson, T.\ 1995, Phys.\ Rev.\ Lett., 75, 1260
\bibitem[Joos \& Zeh(1985)]{JoosZeh1985} Joos, E.\ \& Zeh, H.\ D.\ 1985, Z.\ Phys.\ B, 59, 223
\bibitem[Landauer(1961)]{Landauer1961} Landauer, R.\ 1961, IBM J.\ Res.\ Dev., 5, 183
\bibitem[Mahaffey(2026a)]{IAM_theory} Mahaffey, H.\ W.\ 2026a, Horizon Thermodynamics and Gravitational Decoherence as the Origin of $\mu<1$, $\Sigma=1$, Zenodo, 10.5281/zenodo.18702042
\bibitem[Mahaffey(2026b)]{IAM_obs} Mahaffey, H.\ W.\ 2026b, Dual-Sector Perturbation Cosmology: CAMB Implementation and Validation, Zenodo
\bibitem[Mahaffey(2026c)]{IAM_time} Mahaffey, H.\ W.\ 2026c, The Two Faces of Time: Coordinate Time, Proper Time, and Accumulated Decoherence, Zenodo
\bibitem[Mahaffey(2026d)]{IAM_QM} Mahaffey, H.\ W.\ 2026d, The Measurement Problem Dissolved: Quantum Decoherence as Irreversible Sector Crossing in the Informational Actualization Model, Zenodo
\bibitem[Mahaffey(2026e)]{IAM_grav} Mahaffey, H.\ W.\ 2026e, Gravitational Decoherence from Dual-Sector Thermodynamics: Predictions for Optomechanical Experiments and Quantum Computing Architecture, Zenodo
\bibitem[Mahaffey(2026f)]{IAM_Zurek} Mahaffey, H.\ W.\ 2026f, Quantum Darwinism at Cosmological Scales: Gravitational Decoherence, the Cosmic Horizon, and the Emergence of Classical Structure, Zenodo
\bibitem[Megidish et al.(2013)]{Megidish2013} Megidish, E., et al.\ 2013, Phys.\ Rev.\ Lett., 110, 210403
\bibitem[Scientific Advances(2025)]{SciAdv2025} Bell-inequality violation from path indistinguishability, Sci.\ Adv., 2025
\bibitem[Yin et al.(2017)]{Yin2017} Yin, J., et al.\ 2017, Science, 356, 1140
\bibitem[Zurek(1981)]{Zurek1981} Zurek, W.\ H.\ 1981, Phys.\ Rev.\ D, 24, 1516
\bibitem[Zurek(2003)]{Zurek2003} Zurek, W.\ H.\ 2003, Rev.\ Mod.\ Phys., 75, 715
\bibitem[Zurek(2009)]{Zurek2009} Zurek, W.\ H.\ 2009, Nature Physics, 5, 181
\end{thebibliography}
\end{document}
